\section{圆锥曲线的联立}
\subsection{位置关系}

\begin{theorem}
    椭圆\(C:\dfrac{x^2}{a^2}+\dfrac{y^2}{b^2}=1\)，\(P(x_0.y_0)\)为\(C\)上一点。则过\(P\)点的切线方程为
    \begin{equation}
        \dfrac{x_0x}{a^2}+\dfrac{y_0y}{b^2}=1
    \end{equation}
\end{theorem}
\begin{corollary}
    椭圆\(C:\dfrac{x^2}{a^2}+\dfrac{y^2}{b^2}=1\)，\(P(x_0.y_0)\)为\(C\)外一点。则过\(P\)点对应的切点弦方程为
    \begin{equation}
        \dfrac{x_0x}{a^2}+\dfrac{y_0y}{b^2}=1
    \end{equation}
\end{corollary}




\begin{example}\label{高考:2024年天津卷15:2}
    若函数\(f(x)=2\sqrt{x^2-ax}-|ax-2|+1\)恰有一个零点, 则\(a\)的取值范围\uline{\makebox[1.5cm]{}}. 
\end{example}











\clearpage
\subsection{韦达定理}
\subsubsection{弦长}
\clearpage
\subsubsection{向量关系}
\clearpage
\subsubsection{斜率关系}
\clearpage
\subsubsection{焦半径}

